\chapter{Moments and cumulants}
\label{ch:moments-cumulants}

The paper's laboratory is built on one arithmetic fact: a linear recursion with
independent innovations has cumulants that add. That fact fixes the covariance of
the chain independently of the innovation law, places every non-Gaussian feature in
the cumulants of order four and higher, and gives the fourth-cumulant tensor $\kfour$
that appears in the constants $\tfrac16\Frob{\kfour}^2$ and $\tfrac1{24}\Frob{\partial_\theta\kfour}^2$.
This chapter builds cumulants from moments, derives their three properties (shift,
homogeneity, additivity), extends them to tensors of a random vector, computes the
cumulants of the stationary AR(1) marginal, derives the excess kurtosis of every
innovation family the paper uses, and ends with the exponential-moment assumption
under which the paper's Theorem 2 is proved.

\section{Moments and central moments}

Let $X$ be a real random variable. Its \term{moment} of order $n$ is
$m_n=\E[X^n]$, defined whenever $\E|X|^n<\infty$. Its \term{central moment} of order
$n$ is $\mu_n=\E[(X-m_1)^n]$. If $\E|X|^n<\infty$ then $\E|X|^k<\infty$ for every
$k\le n$, because $|x|^k\le1+|x|^n$; so the moments that exist form an initial segment
of the orders. The variance is $\mu_2=\Var(X)=m_2-m_1^2$.

A law is \term{symmetric} if $-X$ has the law of $X$; then every odd moment that
exists vanishes, $m_{2k+1}=0$, because $\E[X^{2k+1}]=\E[(-X)^{2k+1}]=-\E[X^{2k+1}]$.
Every innovation law in the paper is symmetric, which is why the first non-Gaussian
cumulant is the fourth and not the third.

The \term{kurtosis} of a law with finite fourth moment is $\mu_4/\mu_2^2$ and its
\term{excess kurtosis} is $\mu_4/\mu_2^2-3$. The subtraction of $3$ makes the excess
kurtosis of every Gaussian zero, which the next section explains.

\section{Generating functions and cumulants}

\begin{definition}[Moment- and cumulant-generating functions]
The \term{moment-generating function} of $X$ is $M(s)=\E[e^{sX}]$, $s\in\Reals$, an
extended-real quantity in $(0,\infty]$ with $M(0)=1$. The
\term{cumulant-generating function} is $K(s)=\log M(s)$ wherever $M(s)<\infty$.
\end{definition}

The set where $M$ is finite is an interval containing $0$, because $e^{sx}$ is convex
in $s$. Two cases matter for the paper.

\begin{lemma}[Exponential moment and analyticity]
\label{lem:mgf-analytic}
Suppose $\E[e^{\lambda|X|}]<\infty$ for some $\lambda>0$. Then $M(s)<\infty$ for
$|s|<\lambda$, every moment $m_n$ is finite with $m_n\le n!\,\lambda^{-n}\E[e^{\lambda|X|}]$,
and $M(s)=\sum_{n\ge0}m_ns^n/n!$ for $|s|<\lambda$.
\end{lemma}
\begin{proof}
For $|s|<\lambda$, $e^{sX}\le e^{\lambda|X|}$, so $M(s)<\infty$. From
$e^{\lambda|x|}\ge\lambda^n|x|^n/n!$ (one term of the exponential series) we get
$|x|^n\le n!\lambda^{-n}e^{\lambda|x|}$ and the moment bound. Finally
$\sum_n|s|^n|X|^n/n!=e^{|s||X|}\le e^{\lambda|X|}$ is integrable, so the series
$e^{sX}=\sum_ns^nX^n/n!$ may be integrated term by term by dominated convergence.
\end{proof}

If the exponential moment fails, $M(s)$ may be $+\infty$ for every $s\ne0$ even
though many moments are finite; the Student-$t$ law of \cref{sec:expmoment} is the
paper's example. In that case $K$ is not a function whose Taylor series one can take,
but cumulants of every order for which the moments exist are still defined by the
polynomial relations that follow, read as identities of formal power series.

\begin{definition}[Cumulants]
The \term{cumulant} of order $n$ is $\kappa_n=K^{(n)}(0)$, so that
$K(s)=\sum_{n\ge1}\kappa_ns^n/n!$ near the origin when \cref{lem:mgf-analytic} applies.
In general $\kappa_n$ is the polynomial in $m_1,\dots,m_n$ obtained by expanding
$\log(1+\sum_{k\ge1}m_ks^k/k!)$ formally and reading off the coefficient of $s^n/n!$.
\end{definition}

\paragraph{The first four cumulants.}
Centre $X$ so that $m_1=0$ (the effect of a shift is treated below); then
$M(s)=1+\mu_2s^2/2+\mu_3s^3/6+\mu_4s^4/24+O(s^5)$ and, using
$\log(1+u)=u-u^2/2+O(u^3)$ with $u=O(s^2)$,
\begin{equation}
  K(s)=\frac{\mu_2s^2}{2}+\frac{\mu_3s^3}{6}+\frac{\mu_4s^4}{24}
      -\frac12\Bigl(\frac{\mu_2s^2}{2}\Bigr)^2+O(s^5)
      =\frac{\mu_2}{2}s^2+\frac{\mu_3}{6}s^3+\frac{\mu_4-3\mu_2^2}{24}s^4+O(s^5).
\end{equation}
Reading off coefficients of $s^n/n!$,
\begin{equation}
  \kappa_1=m_1,\qquad \kappa_2=\mu_2,\qquad \kappa_3=\mu_3,\qquad \kappa_4=\mu_4-3\mu_2^2.
  \label{eq:kappa1234}
\end{equation}
So the fourth cumulant is the fourth central moment minus three times the squared
variance, and the excess kurtosis is exactly $\kappa_4/\kappa_2^2$. The factor $3$ is
the number of ways to pair four objects, which \cref{sec:cumulant-tensors} and the
chapter on Gaussian integrals explain.

\begin{proposition}[Three properties of cumulants]
\label{prop:cumulant-props}
Let $X$, $Y$ be random variables with the relevant moments finite, and $c\in\Reals$.
\begin{enumerate}[label=(\roman*)]
\item Shift: $\kappa_1(X+c)=\kappa_1(X)+c$ and $\kappa_n(X+c)=\kappa_n(X)$ for $n\ge2$.
\item Homogeneity: $\kappa_n(cX)=c^n\kappa_n(X)$.
\item Additivity: if $X$ and $Y$ are independent, $\kappa_n(X+Y)=\kappa_n(X)+\kappa_n(Y)$.
\end{enumerate}
\end{proposition}
\begin{proof}
When the generating functions exist near the origin the three statements are read
off $K_{X+c}(s)=cs+K_X(s)$, $K_{cX}(s)=K_X(cs)$ and, by independence,
$M_{X+Y}(s)=M_X(s)M_Y(s)$, hence $K_{X+Y}=K_X+K_Y$. Each statement is a polynomial
identity in finitely many moments, so it holds whenever those moments are finite,
by the formal-series definition.
\end{proof}

\begin{proposition}[Gaussian cumulants]
\label{prop:gauss-cumulants}
If $X\sim\normal{m}{\sigma^2}$ then $K(s)=ms+\tfrac12\sigma^2s^2$, so $\kappa_1=m$,
$\kappa_2=\sigma^2$ and $\kappa_n=0$ for every $n\ge3$.
\end{proposition}
\begin{proof}
For $Z\sim\normal{0}{1}$, completing the square in the exponent,
$\E[e^{sZ}]=(2\pi)^{-1/2}\int e^{sz-z^2/2}\dd z=e^{s^2/2}(2\pi)^{-1/2}\int e^{-(z-s)^2/2}\dd z=e^{s^2/2}$;
then $X=m+\sigma Z$ and \cref{prop:cumulant-props} give the rest.
\end{proof}

A theorem of Marcinkiewicz says the converse: if $K$ is a polynomial then its degree
is at most two and the law is Gaussian. So ``non-Gaussian'' means ``some cumulant of
order at least three is non-zero'', and for a symmetric law ``some cumulant of order
at least four''. This is the sense in which the paper's families ``differ only through
cumulants of order four and higher''.

\why{Additivity and homogeneity are what make the channel act on cumulants
diagonally. If $x_t=e^{-t}a+\sqrt{\Dt}\,z$ with $z$ Gaussian and independent of
$a$, then $\kappa_m(x_t)=e^{-mt}\kappa_m(a)+\Dt^{m/2}\kappa_m(z)=e^{-mt}\kappa_m(a)$ for
$m\ge3$. This is the paper's first ``exact consequence'' of Theorem 1, and it needs
nothing but this section.}

\section{Cumulant tensors of a random vector}
\label{sec:cumulant-tensors}

Let $\vect{X}=(X_1,\dots,X_{\nsites})$ be a random vector. The joint generating
functions are $M(\vect{s})=\E[e^{\vect{s}^{\!\top}\vect{X}}]$ and $K(\vect{s})=\log M(\vect{s})$,
and the \term{joint cumulant} of the coordinates $X_{i_1},\dots,X_{i_n}$ is
\begin{equation}
  \operatorname{cum}(X_{i_1},\dots,X_{i_n})=\frac{\partial^nK}{\partial s_{i_1}\cdots\partial s_{i_n}}(\vect{0}),
\end{equation}
a quantity symmetric in its arguments and multilinear in them. Collected over all
index tuples it forms the \term{cumulant tensor} of order $n$, an array with
$\nsites^n$ entries. The order-two tensor is the covariance matrix, since
$\partial_i\partial_jK(\vect 0)=\E[X_iX_j]-\E[X_i]\E[X_j]$. The order-four tensor is
the object the paper writes $\kfour$.

\begin{lemma}[Fourth joint cumulant of a zero-mean vector]
\label{lem:kappa4-tensor}
If $\E[\vect X]=\vect 0$ and fourth moments are finite, with $\Sigma_{ij}=\E[X_iX_j]$,
\begin{equation}
  \kappa_4[ijkl]=\E[X_iX_jX_kX_l]-\Sigma_{ij}\Sigma_{kl}-\Sigma_{ik}\Sigma_{jl}-\Sigma_{il}\Sigma_{jk}.
  \label{eq:kappa4-tensor}
\end{equation}
\end{lemma}
\begin{proof}
Expand
\[
  M(\vect s)=1+\tfrac12\sum_{ij}\Sigma_{ij}s_is_j+\tfrac16\sum_{ijk}\E[X_iX_jX_k]s_is_js_k
  +\tfrac1{24}\sum_{ijkl}\E[X_iX_jX_kX_l]s_is_js_ks_l+O(|\vect s|^5)
\]
and take $\log(1+u)=u-u^2/2+\dots$; the only fourth-order contribution of $-u^2/2$ is
\[
  -\tfrac18\Bigl(\sum_{ij}\Sigma_{ij}s_is_j\Bigr)^2=-\tfrac18\sum_{ijkl}\Sigma_{ij}\Sigma_{kl}s_is_js_ks_l .
\]
The coefficient of $s_is_js_ks_l$ in $K$, symmetrised, is therefore
$\tfrac1{24}\bigl(\E[X_iX_jX_kX_l]-3\,\mathrm{sym}(\Sigma_{ij}\Sigma_{kl})\bigr)$, where the
symmetrisation of $\Sigma_{ij}\Sigma_{kl}$ over the $24$ permutations of $(i,j,k,l)$
produces each of the three pairings $8$ times, that is
\[
  \mathrm{sym}(\Sigma_{ij}\Sigma_{kl})=\tfrac13\bigl(\Sigma_{ij}\Sigma_{kl}+\Sigma_{ik}\Sigma_{jl}+\Sigma_{il}\Sigma_{jk}\bigr).
\]
Multiplying by $4!$ gives \cref{eq:kappa4-tensor}. On the diagonal $i=j=k=l$ it
reduces to \cref{eq:kappa1234}.
\end{proof}

The three properties of \cref{prop:cumulant-props} have tensor versions, proved the
same way from $K$. Shifts change only the first-order tensor. Independence gives
additivity: if $\vect X$ and $\vect Y$ are independent then the cumulant tensors of
$\vect X+\vect Y$ are the sums of those of $\vect X$ and of $\vect Y$. Homogeneity
becomes the \term{transformation rule} for a linear map $\mat A$:
\begin{equation}
  \operatorname{cum}\bigl((\mat A\vect X)_{i_1},\dots,(\mat A\vect X)_{i_n}\bigr)
  =\sum_{j_1,\dots,j_n}A_{i_1j_1}\cdots A_{i_nj_n}\operatorname{cum}(X_{j_1},\dots,X_{j_n}),
  \label{eq:cum-transform}
\end{equation}
which follows from $K_{\mat A\vect X}(\vect s)=K_{\vect X}(\mat A^{\!\top}\vect s)$ and the
chain rule. A Gaussian vector has $K(\vect s)=\vect m^{\!\top}\vect s+\tfrac12\vect s^{\!\top}\Sig\vect s$
and no cumulant tensor of order three or more.

\paragraph{The Frobenius norm over ordered indices.}
For a tensor $\mat T$ of order $r$, $\Frob{\mat T}^2=\sum_{i_1\cdots i_r}T_{i_1\cdots i_r}^2$,
the sum over all ordered tuples. This is the Euclidean norm of $\mat T$ viewed as a
vector with $\nsites^r$ entries, and it is invariant under orthogonal changes of
coordinates, because by \cref{eq:cum-transform} an orthogonal $\mat O$ acts on the
tensor by $\mat O\otimes\cdots\otimes\mat O$, which is orthogonal on $\Reals^{\nsites^r}$.
For a symmetric tensor the ordered sum can be regrouped by \term{multisets}: an ordered
tuple with multiplicities $\nu=(\nu_1,\dots,\nu_{\nsites})$, $|\nu|=r$, has $r!/\nu!$
rearrangements, all carrying the same entry $T_\nu$, so
\begin{equation}
  \Frob{\mat T}^2=\sum_{|\nu|=r}\frac{r!}{\nu!}\,T_\nu^2,
  \qquad\text{equivalently}\qquad
  \sum_{|\nu|=r}\frac{T_\nu^2}{\nu!}=\frac{1}{r!}\Frob{\mat T}^2.
  \label{eq:frob-multiset}
\end{equation}
For $r=4$ the right-hand side is $\Frob{\mat T}^2/24$. This is the origin of the paper's
constant $\tfrac1{24}\Frob{\partial_\theta\kfour}^2$: the general-order theorem gives
$\sum_{|\nu|=4}(\partial_\theta\E H_\nu)^2/\nu!$, the degree-four Hermite moments of a
variance-matched family are the entries of $\kfour$ (chapter on Gaussian integrals and
Hermite polynomials), and \cref{eq:frob-multiset} converts the multiset sum into the
ordered-index norm. For a symmetric matrix ($r=2$), $\sum_{i\le j}$ with weights
$2$ off the diagonal gives $\tfrac12\Frob{\partial_\theta\Sig_0}^2=\sum_i\tfrac12(\partial\Sigma_{ii})^2+\sum_{i<j}(\partial\Sigma_{ij})^2$,
the paper's constant for a covariance coordinate.

\section{Cumulants of the stationary AR(1) marginal}
\label{sec:ar1-cumulants}

Let $(\varepsilon_k)_{k\in\mathbb Z}$ be independent copies of a law with mean zero,
variance $\ivar$ and finite cumulants up to the order considered, and let $|\corr|<1$.
The \term{moving-average representation}
\begin{equation}
  X=\sum_{k\ge0}\corr^k\varepsilon_{-k}
  \label{eq:ma-rep}
\end{equation}
converges in mean square (its partial sums are Cauchy in $L^2$ because
$\sum_k\corr^{2k}<\infty$) and almost surely, and $X$ has the \term{invariant law} of
the recursion $a_i=\corr a_{i-1}+\varepsilon_i$: if $a_{i-1}$ has the law of $X$ then
so does $a_i$, since $\corr X+\varepsilon$ is again of the form \cref{eq:ma-rep}. The
chapter on Markov chains returns to this.

\begin{proposition}[Cumulants of the AR(1) marginal]
\label{prop:ar1-cumulants}
If $\E|\varepsilon|^n<\infty$ then $\E|X|^n<\infty$ and
\begin{equation}
  \kappa_n(X)=\frac{\kappa_n(\varepsilon)}{1-\corr^n},\qquad n\ge1.
  \label{eq:ar1-cumulants}
\end{equation}
\end{proposition}
\begin{proof}
For the finite sum $X_N=\sum_{k=0}^N\corr^k\varepsilon_{-k}$, additivity and homogeneity
give $\kappa_n(X_N)=\kappa_n(\varepsilon)\sum_{k=0}^N\corr^{nk}$. Since
$\E|X_N-X|^n\to0$ (Minkowski's inequality and the tail $\sum_{k>N}|\corr|^k$), every
moment of order at most $n$ of $X_N$ converges to that of $X$, so the cumulants,
polynomials in those moments, converge too. Summing the geometric series gives
\cref{eq:ar1-cumulants}.
\end{proof}

Two consequences fix the paper's laboratory. With $\ivar=1-\corr^2$,
\begin{equation}
  \Var(X)=\frac{\ivar}{1-\corr^2}=1,
\end{equation}
so the choice of innovation variance is exactly the one that gives unit marginal
variance, whatever the innovation law. And the fourth cumulant and excess kurtosis
of the marginal are
\begin{equation}
  \kappa_4(X)=\frac{\kappa_4(\varepsilon)}{1-\corr^4},\qquad
  \frac{\kappa_4(X)}{\Var(X)^2}=\frac{\kappa_4(\varepsilon)}{\ivar^2}\cdot\frac{\ivar^2}{1-\corr^4}
  =\gamma_2(\varepsilon)\,\frac{(1-\corr^2)^2}{(1-\corr^2)(1+\corr^2)}
  =\gamma_2(\varepsilon)\,\frac{1-\corr^2}{1+\corr^2},
  \label{eq:ar1-kurtosis}
\end{equation}
where $\gamma_2(\varepsilon)=\kappa_4(\varepsilon)/\ivar^2$ is the innovation's excess
kurtosis. For Laplace innovations, $\gamma_2(\varepsilon)=3$ (next section), the
marginal's excess kurtosis is $3(1-\corr^2)/(1+\corr^2)$, which is the expression the
paper uses in its regime sweep (Section 5) and in the hypothesis of E14. It falls from
$3$ at $\corr=0$ to $0$ as $\corr\to1$: a strongly correlated chain has a nearly
Gaussian marginal, because each site is the sum of many independent innovations with
slowly decaying weights, and cumulants of order four add with weights $\corr^{4k}$
while the variance adds with weights $\corr^{2k}$.

\begin{remark}[The paper's chain is not started in its invariant law]
The paper's chain has $a_1\sim\normal{0}{1}$, which is the invariant law only for
Gaussian innovations. For every other family the first site has no fourth cumulant,
and the site-$i$ marginal has
\[
  \kappa_4(a_i)=\kappa_4(\varepsilon)\sum_{m=2}^{i}\corr^{4(i-m)}
  =\kappa_4(\varepsilon)\,\frac{1-\corr^{4(i-1)}}{1-\corr^4},
\]
which rises geometrically to the stationary value \cref{eq:ar1-kurtosis}. The
covariance is unaffected (chapter on Markov chains); the fourth-cumulant tensor of
\code{src/markov\_score/theory.py} (function \code{chain\_kappa4}) carries the exact
site-dependent form, and the moment estimator of E19
(\code{experiments/e19\_moment\_estimators.py}) uses the site-dependent factor
$(1-\corr^{4(i-1)})/(1-\corr^4)$ when it converts a per-site fourth cumulant into an
estimate of $\kappa_4(\varepsilon)$. The paper's sentence
``its excess kurtosis falling as $3(1-\corr^2)/(1+\corr^2)$'' refers to the stationary
marginal.
\end{remark}

\section{The innovation families at variance $\ivar$}
\label{sec:families}

Every family below is symmetric, so $\kappa_3=0$, and is parameterised so that its
variance is exactly $\ivar$. The remaining freedom, if any, is a shape parameter. For
each we derive the fourth moment by hand and record the excess kurtosis
$\gamma_2=\mu_4/\ivar^2-3$ and the fourth cumulant $\kappa_4=\gamma_2\ivar^2$. The
parameterisations are those of \code{src/markov\_score/priors.py} and
\code{src/markov\_score/shape\_kernel.py}.

\paragraph{Gaussian.}
$\varepsilon\sim\normal{0}{\ivar}$. By \cref{prop:gauss-cumulants}, $\kappa_4=0$ and
$\gamma_2=0$; $\mu_4=3\ivar^2$.

\paragraph{Laplace.}
Density $\innov(e)=\tfrac{1}{2b}e^{-|e|/b}$. The absolute moments are
$\E|\varepsilon|^n=\tfrac1b\int_0^\infty e^ne^{-e/b}\dd e=n!\,b^n$ (the Gamma integral),
so $\mu_2=2b^2$ and $\mu_4=24b^4$. Variance matching sets $2b^2=\ivar$, i.e.
$b=\sqrt{\ivar/2}$, and then
\begin{equation}
  \gamma_2=\frac{24b^4}{(2b^2)^2}-3=6-3=3,\qquad \kappa_4=3\ivar^2.
\end{equation}

\paragraph{Student-$t_\nu$.}
$\varepsilon=sT$ with $T$ a Student-$t$ variable with $\nu$ degrees of freedom,
density proportional to $(1+u^2/\nu)^{-(\nu+1)/2}$. Writing $T=Z/\sqrt{V/\nu}$ with
$Z\sim\normal01$ and $V\sim\chi^2_\nu$ independent, $\E[T^{2k}]=\E[Z^{2k}]\,\nu^k\E[V^{-k}]$,
and $\E[V^{-k}]=\prod_{j=1}^k(\nu-2j)^{-1}$ for $\nu>2k$ (the inverse moments of a
chi-square). Hence $\E[T^2]=\nu/(\nu-2)$ for $\nu>2$ and $\E[T^4]=3\nu^2/((\nu-2)(\nu-4))$
for $\nu>4$. Variance matching sets $s^2\nu/(\nu-2)=\ivar$, and
\begin{equation}
  \gamma_2=\frac{3\nu^2}{(\nu-2)(\nu-4)}\cdot\frac{(\nu-2)^2}{\nu^2}-3
  =\frac{3(\nu-2)}{\nu-4}-3=\frac{6}{\nu-4},\qquad \kappa_4=\frac{6}{\nu-4}\ivar^2.
\end{equation}
At $\nu=5$, the paper's choice, $\gamma_2=6$. Moments of order $r$ exist only for
$r<\nu$: the $t_5$ law has finite moments up to order four and an infinite sixth moment,
so its cumulants beyond the fourth do not exist. Its tail is polynomial,
$\innov(e)\sim c|e|^{-\nu-1}$, which is what \cref{sec:expmoment} is about.

\paragraph{Uniform.}
$\varepsilon\sim\mathrm{Uniform}(-h,h)$, density $1/(2h)$ on $[-h,h]$. Then
$\mu_2=\tfrac1{2h}\int_{-h}^he^2\dd e=h^2/3$ and $\mu_4=h^4/5$. Variance matching sets
$h=\sqrt{3\ivar}$, and
\begin{equation}
  \gamma_2=\frac{h^4/5}{h^4/9}-3=\frac95-3=-\frac65=-1.2,\qquad \kappa_4=-\frac65\ivar^2.
\end{equation}
This is the most platykurtic law of the paper, and it is compactly supported: the
conditional law of $a_i$ given $a_{i-1}$ is zero outside an interval, which is a
different failure of Gaussian closure from heavy tails.

\paragraph{Symmetric two-component mixture.}
$\varepsilon\sim\tfrac12\normal{m}{s^2}+\tfrac12\normal{-m}{s^2}$ with $m^2=\kappa\ivar$
and $s^2=(1-\kappa)\ivar$, $\kappa\in[0,1)$; $\kappa$ is the separation fraction and
the frozen value is $\kappa=0.6$. For the $+m$ component, $\E[(m+sZ)^4]=m^4+6m^2s^2+3s^4$
(binomial expansion and Gaussian moments $\E Z^2=1$, $\E Z^4=3$, odd ones zero); the
$-m$ component gives the same, so $\mu_4=m^4+6m^2s^2+3s^4$ and $\mu_2=m^2+s^2=\ivar$.
Substituting,
\begin{equation}
  \gamma_2=\frac{\kappa^2+6\kappa(1-\kappa)+3(1-\kappa)^2}{1}-3
  =\kappa^2+6\kappa-6\kappa^2+3-6\kappa+3\kappa^2-3=-2\kappa^2,\qquad \kappa_4=-2\kappa^2\ivar^2.
\end{equation}
At $\kappa=0.6$, $\gamma_2=-0.72$, the value in the paper's oracle table. The law is
bimodal for $\kappa$ large enough, and at $\kappa=0$ it is the Gaussian.

\paragraph{Generalised Gaussian with exponent $\beta$.}
Density $\innov_\beta(e)=\dfrac{\beta}{2s\,\Gamma(1/\beta)}\exp\bigl(-(|e|/s)^\beta\bigr)$,
$\beta>0$. The substitution $u=(e/s)^\beta$, $e=su^{1/\beta}$, $\dd e=(s/\beta)u^{1/\beta-1}\dd u$
turns the absolute moments into Gamma integrals:
\begin{equation}
  \E|\varepsilon|^k=\frac{\beta}{s\,\Gamma(1/\beta)}\int_0^\infty e^ke^{-(e/s)^\beta}\dd e
  =\frac{s^k}{\Gamma(1/\beta)}\int_0^\infty u^{(k+1)/\beta-1}e^{-u}\dd u
  =s^k\,\frac{\Gamma\bigl((k+1)/\beta\bigr)}{\Gamma(1/\beta)}.
\end{equation}
(At $k=0$ this confirms the normalisation.) Hence $\mu_2=s^2\Gamma(3/\beta)/\Gamma(1/\beta)$,
so variance matching sets $s_\beta=\sqrt{\ivar\,\Gamma(1/\beta)/\Gamma(3/\beta)}$, and
\begin{equation}
  \gamma_2(\beta)=\frac{\Gamma(5/\beta)\,\Gamma(1/\beta)}{\Gamma(3/\beta)^2}-3.
  \label{eq:gg-kurtosis}
\end{equation}
Three values by hand. At $\beta=1$: $\Gamma(5)\Gamma(1)/\Gamma(3)^2-3=24/4-3=3$, the
Laplace law (indeed $s_1=\sqrt{\ivar\cdot1/2}=b$). At $\beta=2$:
$\Gamma(5/2)\Gamma(1/2)/\Gamma(3/2)^2-3=(\tfrac34\sqrt\pi)(\sqrt\pi)/(\tfrac12\sqrt\pi)^2-3=3-3=0$,
the Gaussian. At $\beta=6$: $\Gamma(5/6)\Gamma(1/6)/\Gamma(1/2)^2-3$, and by the
reflection formula $\Gamma(z)\Gamma(1-z)=\pi/\sin(\pi z)$, $\Gamma(1/6)\Gamma(5/6)=\pi/\sin(\pi/6)=2\pi$
while $\Gamma(1/2)^2=\pi$; so $\gamma_2(6)=2-3=-1$ exactly. As $\beta\to\infty$ the law
tends to the uniform and $\gamma_2\to-6/5$. The exponent $\beta$ therefore moves the
excess kurtosis continuously from $+3$ through $0$ to $-1$ on the range $\beta\in[1,6]$
that the paper's protocol declares for E02b, at fixed variance and (for the chain)
fixed covariance; the paper's Section 2 quotes the platykurtic end as $-1.1$, a value
that \cref{eq:gg-kurtosis} reaches only near $\beta\approx9.5$. Along the family
$\ivar$ is fixed and $\beta$ is a pure \term{shape coordinate}: $\partial_\beta\Sig_0=0$
and $\partial_\beta\kfour\ne0$, which is exactly what the paper's Theorem 2 needs for
its $e^{-8t}$ statement.

\begin{table}[t]
\centering
\caption{The innovation families at variance $\ivar$: parameter, excess kurtosis
$\gamma_2$ (the fourth cumulant is $\kappa_4=\gamma_2\ivar^2$), and the range of
$\lambda$ for which the exponential moment of \cref{sec:expmoment} is finite.}
\label{tab:families}
\footnotesize
\begin{tabular}{@{}llcl@{}}
\toprule
family & parameter at variance $\ivar$ & $\gamma_2$ & $\E e^{\lambda|\varepsilon|}<\infty$ \\
\midrule
Gaussian & --- & $0$ & all $\lambda$ \\
Laplace & $b=\sqrt{\ivar/2}$ & $3$ & $\lambda<1/b$ \\
Student-$t_\nu$ & $s=\sqrt{\ivar(\nu-2)/\nu}$ & $6/(\nu-4)$ & no $\lambda$ \\
uniform & $h=\sqrt{3\ivar}$ & $-6/5$ & all $\lambda$ \\
mixture, separation $\kappa$ & $m^2=\kappa\ivar$, $s^2=(1-\kappa)\ivar$ & $-2\kappa^2$ & all $\lambda$ \\
generalised Gaussian, $\beta$ & $s_\beta=\sqrt{\ivar\Gamma(1/\beta)/\Gamma(3/\beta)}$ & \cref{eq:gg-kurtosis} & all if $\beta>1$; $\lambda<1/s_1$ if $\beta=1$ \\
\bottomrule
\end{tabular}
\end{table}

\section{The exponential-moment assumption}
\label{sec:expmoment}

\begin{assumption}[Exponential moment]
\label{ass:expmoment}
A law has an \term{exponential moment} if $\E[e^{\lambda|X|}]<\infty$ for some $\lambda>0$.
For a random vector the condition is $\E[e^{\lambda|\vect X|}]<\infty$ with $|\cdot|$
the Euclidean norm.
\end{assumption}

This is the paper's Assumption 1, the one hypothesis under which its Theorem 2 is
proved. Its consequences, all used in that proof, are collected here.

\begin{proposition}
\label{prop:expmoment-consequences}
Under \cref{ass:expmoment}:
\begin{enumerate}[label=(\roman*)]
\item all moments are finite, and the cumulant-generating function $K(\vect s)$ is
  finite and real-analytic on $\{|\vect s|<\lambda\}$;
\item the law is determined by its moments;
\item for the scalar Student-$t_\nu$ law the assumption fails for every $\lambda>0$,
  and $K(s)=+\infty$ for every $s\ne0$.
\end{enumerate}
\end{proposition}
\begin{proof}
(i) is \cref{lem:mgf-analytic} applied to $\vect s^{\!\top}\vect X\le|\vect s||\vect X|$; the
power series of $M$ converges on the ball and $\log$ of a positive analytic function
is analytic. (ii) The characteristic function $\phi(\vect s)=\E[e^{\mathrm i\vect s^{\!\top}\vect X}]$
extends to an analytic function of complex $\vect s$ on the strip $|\operatorname{Im}\vect s|<\lambda$,
because $\E[e^{\lambda|\vect X|}]<\infty$ dominates the integrand there; an analytic
function on a strip containing the real axis is determined by its derivatives at the
origin, which are the moments up to powers of $\mathrm i$; and the characteristic
function determines the law. (iii) The density of $T$ satisfies $\innov(u)\ge c\,u^{-\nu-1}$
for $u\ge u_0$, so for $s>0$,
$\E[e^{sT}]\ge c\int_{u_0}^\infty e^{su}u^{-\nu-1}\dd u=\infty$, because
$e^{su}u^{-\nu-1}\to\infty$; by symmetry the same holds for $s<0$. Then $M(s)=\infty$
for all $s\ne0$, so $\E[e^{\lambda|T|}]\ge\E[e^{\lambda T}]=\infty$.
\end{proof}

\paragraph{What fails without it.}
For the Student-$t$ law the moments up to order $\lfloor\nu\rfloor$ exist and so do
the cumulants up to that order, as polynomials in them; but $K$ is not the sum of the
series $\sum\kappa_ns^n/n!$, because $K$ is infinite off the origin. A Taylor expansion
of $K$ with a Lagrange remainder therefore does not exist, and this is the reason the
paper's Theorem 2 excludes Student-$t$ and treats it separately as a heavy-tailed
regime with a rate of its own, supported by quadrature and not claimed as a theorem.
The paper's damped generating function
$\Lambda(\vect s,b)=\log\E_{\Pz}\exp(\vect s^{\!\top}\va-b|\va|^2/2)$ is the cumulant-generating
function of the linear form $\vect s^{\!\top}\va$, damped by a Gaussian factor; for $b>0$
the damping makes it finite for every $\vect s$ and every law, but the expansion the
theorem needs is at $b=0$, where $\Lambda(\vect s,0)=K(\vect s)$, and \cref{ass:expmoment}
is what makes that expansion legitimate on $|\vect s|<\lambda$ uniformly in $b\ge0$.

\paragraph{Which families satisfy it.}
The Gaussian, the uniform and the mixture have all exponential moments (the Gaussian
by $\E e^{\lambda|Z|}\le2e^{\lambda^2/2}$, the others by bounded support or by being
mixtures of Gaussians). The Laplace law with scale $b$ has
$\E e^{\lambda|\varepsilon|}=\tfrac1b\int_0^\infty e^{(\lambda-1/b)e}\dd e=(1-\lambda b)^{-1}$
for $\lambda<1/b$ and $+\infty$ for $\lambda\ge1/b$: the assumption holds, with a
finite range of $\lambda$. The generalised Gaussian with $\beta>1$ has all exponential
moments, since $\lambda e-(e/s)^\beta\to-\infty$; at $\beta=1$ it is the Laplace case.
Student-$t$ fails, by (iii). This is the last column of \cref{tab:families}.

\paragraph{From the innovation to the chain.}
The chain $\va$ has an exponential moment if and only if its innovation does. If it
does, write $a_i=\corr^{i-1}a_1+\sum_{m=2}^i\corr^{i-m}\varepsilon_m$; then
$|\va|\le\sum_i|a_i|\le C|a_1|+C\sum_m|\varepsilon_m|$ with $C=\sum_{k\ge0}|\corr|^k=(1-|\corr|)^{-1}$,
so by independence
\[
  \E e^{\lambda|\va|}\le\E e^{\lambda C|a_1|}\prod_{m=2}^{\nsites}\E e^{\lambda C|\varepsilon_m|}<\infty
\]
for $\lambda C$ below the innovation's threshold. Conversely, if the chain has one
then so does $a_2$, and from $|\varepsilon_2|\le|a_2|+|\corr||a_1|$ and the
Cauchy--Schwarz inequality,
\[
  \E e^{\lambda|\varepsilon_2|}\le\bigl(\E e^{2\lambda|a_2|}\bigr)^{1/2}\bigl(\E e^{2\lambda|\corr||a_1|}\bigr)^{1/2}<\infty
\]
for $\lambda$ small enough. So the assumption is a property of the innovation family,
and \cref{tab:families} settles it for the chain as well.

\section{Worked example: the Laplace chain at $\corr=0.85$}

Take the frozen configuration. The innovation variance is
$\ivar=1-\corr^2=1-0.7225=0.2775$, so the Laplace scale is
$b=\sqrt{\ivar/2}=\sqrt{0.13875}\approx0.3725$. The innovation's fourth cumulant is
$\kappa_4(\varepsilon)=3\ivar^2=3\times0.07701\approx0.2310$, and its excess kurtosis
is $3$ whatever $\corr$ is.

The stationary marginal has, by \cref{eq:ar1-cumulants}, $\kappa_4(X)=\kappa_4(\varepsilon)/(1-\corr^4)$
with $\corr^4=0.7225^2\approx0.5220$, so $\kappa_4(X)\approx0.2310/0.4780\approx0.4833$;
since $\Var(X)=1$ this is also the marginal excess kurtosis, and \cref{eq:ar1-kurtosis}
gives the same number directly: $3\times0.2775/1.7225\approx0.4833$. The marginal is
six times less kurtotic than the innovation: most of the fourth-order structure of the
innovation is averaged away by the correlation, which is the direction of the paper's
regime sweep.

Along the paper's chain, which starts at a Gaussian site, the site-$i$ fourth cumulant
is $\kappa_4(\varepsilon)(1-\corr^{4(i-1)})/(1-\corr^4)$: it is $0$ at site $1$,
$\kappa_4(\varepsilon)\approx0.231$ at site $2$, $\kappa_4(\varepsilon)(1+\corr^4)\approx0.352$
at site $3$, $\approx0.415$ at site $4$, and within one per cent of the stationary
value from site $9$ on, since $\corr^{32}\approx0.0056$. The exponential moment holds
for $\lambda<1/b\approx2.68$.

The scalar leading constant of the paper's closure gap, $\kappa_4^2/6$ for a unit-variance
law, is $9/6=1.5$ for the unit-variance Laplace law (excess kurtosis $3$, so
$\kappa_4=3$); the chain constant $\tfrac16\Frob{\kfour}^2$ is a sum of squares of the
site-dependent tensor of the chapter on Markov chains, and its numerical value at
$\nsites=32$ is a generated macro of the paper, not typed here.

\where{The paper's Section 2 (variance matching, the families and their excess
kurtoses), Section 3 (the two exact consequences of Theorem 1, the constants
$\tfrac16\Frob{\kfour}^2$ and $\tfrac1{24}\Frob{\partial_\theta\kfour}^2$, Assumption 1
and the heavy-tailed paragraph), Section 5 (the marginal excess kurtosis in the regime
sweep) and Appendix A (the grading of $\Lambda$ by cumulants). Experiments E01, E02b,
E13, E14, E18, E19, E21. Code: \code{src/markov\_score/priors.py} (the families and
their \code{innovation\_excess\_kurtosis} properties),
\code{src/markov\_score/shape\_kernel.py} (the generalised Gaussian, its scale at fixed
variance and \code{generalized\_gaussian\_excess\_kurtosis}),
\code{src/markov\_score/theory.py} (\code{chain\_kappa4}, \code{score\_gap\_constant},
\code{excess\_kurtosis\_to\_kappa4}), \code{tests/test\_theory\_hermite.py}
(the fourth-cumulant tensor against brute-force cumulants of sampled chains).}

\paragraph{References.}
P. McCullagh (1987), \emph{Tensor Methods in Statistics}, Chapman and Hall (Chapters 2
and 3 for cumulant tensors and the transformation rule).
P. Billingsley (1995), \emph{Probability and Measure}, 3rd edition, Wiley (Sections 21
and 26 for moments, generating functions and the moment problem).
E. Lukacs (1970), \emph{Characteristic Functions}, 2nd edition, Griffin (analytic
characteristic functions; Marcinkiewicz's theorem).
P. J. Brockwell and R. A. Davis (1991), \emph{Time Series: Theory and Methods},
2nd edition, Springer (the moving-average representation of AR(1)).
S. Nadarajah (2005), ``A generalized normal distribution'', Journal of Applied
Statistics 32, 685--694 (moments of the generalised Gaussian).
